% =====================================================================
%  core/packages.tex — preambuła: pakiety, fonty, ustawienia globalne
%
%  Komentarz "WYMOGI" oznacza wartości narzucone przez wymagania
%  edytorskie wydziału — nie zmieniaj ich bez konsultacji z promotorem.
%  Kolejność ładowania też nie jest dowolna (np. unicode-math PO
%  fontspec, hyperref prawie na końcu) — przy przestawianiu czytaj
%  komentarze.
% =====================================================================

% --- 1. Podstawowa konfiguracja i język ---
\usepackage[a4paper, top=2.5cm, bottom=2.5cm, left=3.5cm, right=2cm]{geometry}
% WYMOGI: marginesy. Lewy 3,5 cm jest większy, bo tam strona schodzi pod oprawę.
\usepackage{indentfirst}
% Wcięcie także w pierwszym akapicie po nagłówku — polska konwencja typograficzna.
\usepackage[polish]{babel}
% Polskie napisy ("Rysunek", "Tabela", "Spis treści"...) i dzielenie wyrazów.
% W XeLaTeX używamy babel, a NIE pakietu 'polski' — tamten jest dla pdfLaTeX.

% --- 2. Fonty i matematyka (XeLaTeX) ---
% fontspec/unicode-math działają tylko pod XeLaTeX/LuaLaTeX — stąd -pdfxe w compile.sh.
\usepackage{mathtools} % nadzbiór amsmath (środowiska matematyczne)
\usepackage{amssymb}
\usepackage{amsthm}
\usepackage{fontspec}
\usepackage{unicode-math} % to musi być PO fontspec i amsmath

% Domyślny krój: Latin Modern. Teoretycznie powinien to być Times New Roman, jednak na
% linuxie jego instalcja jest problematyczna, a LM jest bardzo podobny i w pełni darmowy;
\setmainfont{Latin Modern Roman}
\setsansfont{Latin Modern Sans}

% Mono (kod źródłowy): Latin Modern Mono nie ma znaków indeksów dolnych (₀₁),
% więc bierzemy DejaVu Sans Mono — pełne pokrycie znaków, standard w Linuksie.
% Nie masz DejaVu? Tu wchodzi np. "FreeMono" albo "Courier New".
\setmonofont{DejaVu Sans Mono}[
	Scale=MatchLowercase,
	%FakeBold=1.0, % emulacja pogrubienia dla kodu
]

\setmathfont{Latin Modern Math} % wymagane przy unicode-math

% --- 3. Grafika i kolory ---
\usepackage{graphicx}          % \includegraphics — obrazki PNG/JPG/PDF
\usepackage{wrapfig}           % tekst opływający rysunek
\usepackage[section]{placeins} % rysunki nie uciekają poza swoją sekcję
\usepackage{float}             % opcja [H] = "dokładnie tutaj, nie kombinuj"
\usepackage{tikz}              % rysunki wektorowe w kodzie (przykład: 50-impl.tex)
\usepackage{pgfplots}          % wykresy w TikZ-u
\pgfplotsset{compat=1.18}

% \usepackage[inkscapeformat=png]{svg}
% ^ Odkomentuj, jeśli chcesz wstawiać .svg. Wymaga zainstalowanego Inkscape
%   (konwersja dzieje się przy kompilacji). Bez .svg w pracy — nie włączaj.
%   użyj wtedy \includesvg{pathdopliku}

% --- 4. Tabele i listy ---
\usepackage{paralist}  % kompaktowe listy (compactitem i spółka)
\usepackage{enumerate} % własne etykiety w \begin{enumerate}[...]
\usepackage{multirow}  % komórka na kilka wierszy
\usepackage{multicol}  % kilka łamów w obrębie strony
\usepackage{tabulary}  % tabele z zawijaniem tekstu do szerokości
\usepackage{ragged2e}  % \justify i spółka (używane na stronie tytułowej)
\usepackage{makecell}  % wielolinijkowe komórki \makecell{...}
\usepackage{tocloft}   % styl spisu treści (szerokość kolumny numerów itd.)

% --- 5. Kod i algorytmy ---
\usepackage{minted}
% minted koloruje składnię przez ZEWNĘTRZNY program pygmentize, dlatego:
%  a) musi być zainstalowany (pip install Pygments; test: pygmentize -V),
%  b) kompilacja musi iść z -shell-escape — compile.sh już to robi.
\usepackage[ruled, vlined]{algorithm2e} % pseudokod (przykład: 50-impl.tex)

% --- 6. Bibliografia (biblatex) ---
\usepackage[nottoc]{tocbibind}
% Spisy i bibliografia wchodzą do spisu treści
\usepackage[
	sorting=none,       % numery [1], [2]... w kolejności CYTOWANIA, nie alfabetycznie
	style=numeric-comp, % styl numerowany, [3–6] zamiast [3,4,5,6]
	maxbibnames=3,      % przy 4+ autorach w bibliografii pojawia się "i in."
	minbibnames=3
]{biblatex}
\addbibresource{bibl.bib} % baza źródeł; klucze cytujesz przez \cite{klucz}
\DefineBibliographyStrings{polish}{
	andothers = {i~in\adddot} % polskie "i in." zamiast "et al."
}

% --- 7. Różne ---
\usepackage{siunitx} % jednostki i liczby: \qty{9.81}{m/s^2} -> "9,81 m/s²"
\linespread{1.5}     % WYMOGI: interlinia 1,5
\parindent=1cm       % WYMOGI: wcięcie akapitowe
\setlength{\cftfignumwidth}{4em} % tocloft: miejsce na numery w spisie treści

% Numeracja wewnątrz rozdziałów: wzór (3.1), rysunek 3.1 itd.
\numberwithin{equation}{section}
\numberwithin{figure}{section}
\counterwithin{figure}{section}

% Polskie nazwy algorytmów ("Algorytm 1"; spis: "Spis algorytmów")
\SetAlgorithmName{Algorytm}{algorytm}{Spis algorytmów}

% Biblioteki TikZ. Dokładanie "na zapas" spowalnia kompilację —
% dopisuj, gdy faktycznie czegoś brakuje.
\usetikzlibrary{
	automata, shapes.multipart, positioning, arrows.meta,
	shadows, fit, shapes.geometric, calc, patterns
}

% --- 8. Pakiety ingerujące w PDF (ładować na końcu!) ---
\usepackage{pdfpages} % \includepdf — np. wklejenie karty dyplomowej PDF (0-head.tex)

\usepackage[colorlinks, plainpages=false, pdfpagelabels, naturalnames=true]{hyperref}
\hypersetup{
	linkcolor={black!50!black}, % odnośniki wewnętrzne
	citecolor={blue!50!black}, % cytowania: przygaszony niebieski
	urlcolor={blue!80!black}   % adresy www: niebieski
}

% --- Koniec preambuły ---
